\chapter{Анализ результатов}\label{ch:7}

\section{Результаты обучения моделей компьютерного зрения}

Для решения задачи автоматизированного анализа фотографии на предмет повреждений автомобиля был реализован двухэтапный ML-конвейер. Подход разделения задачи на два этапа - поиска деталей и поиска повреждений - позволил устранить проблему логически некорректных сочетаний (например, <фара - спущенная шина>), а также приблизить модель к реальному сценарию работы автосервиса, где модель, как и человек, анализирует повреждения в рамках отдельных деталей автомобиля.

В ходе работы были обучены две модели компьютерного зрения на базе YOLOv8-seg:

\begin{itemize}
    \item Модель сегментации деталей автомобиля
    \item Модель сегментации повреждений автомобиля
\end{itemize}


\begin{table}[H]
\centering
\caption{Итоговые метрики модели сегментации деталей}
\begin{tabular}{|l|l|}
\hline
\textbf{Метрика} & \textbf{Значение} \\
\hline
Размер обучающей выборки     & 174 изображения \\
\hline
Размер валидационной выборки & 26 изображений \\
\hline
Размер тестовой выборки      & 26 изображений \\
\hline
Precision (Box)              & 0.766 \\
\hline
Recall (Box)                 & 0.517 \\
\hline
mAP@0.5 (Box)                & 0.560 \\
\hline
mAP@0.5:0.95 (Box)           & 0.471 \\
\hline
Начальный box\_loss          & 1.171 \\
\hline
Конечный box\_loss           & 0.3387 \\
\hline
Количество итераций (эпох)   & 150 \\
\hline
\% падения loss              & 346\% \\
\hline
\end{tabular}
\end{table}

\begin{table}[H]
\centering
\caption{Итоговые метрики модели сегментации повреждений}
\begin{tabular}{|l|l|}
\hline
\textbf{Метрика} & \textbf{Значение} \\
\hline
Размер обучающей выборки     & 313 изображений \\
\hline
Размер валидационной выборки & 29 изображений \\
\hline
Размер тестовой выборки      & 30 изображений \\
\hline
Precision (Mask)             & 0.684 \\
\hline
Recall (Mask)                & 0.449 \\
\hline
mAP@0.5 (Mask)               & 0.474 \\
\hline
mAP@0.5:0.95 (Mask)          & 0.307 \\
\hline
Начальный seg\_loss          & 3.463 \\
\hline
Конечный seg\_loss           & 1.377 \\
\hline
Количество итераций          & 137 (137/150 EarlyStopping) \\
\hline
\% падения loss              & 251\% \\
\hline
\end{tabular}
\end{table}

Результаты обучения обеих моделей визуализированы и представлены на графиках.
На рисунках 22 и 23 представлены графики loss и метрик, а на рисунках 24 и 25 - скриншоты процессов обучения моделей.


\begin{figure}[H]
\centering
\includegraphics[scale=0.8]{details_training_combined}
\caption{График метрик обучения модели деталей}
\end{figure}

\begin{figure}[H]
\centering
\includegraphics[scale=0.8]{damages_training_combined}
\caption{График метрик обучения модели повреждений}
\end{figure}

\begin{figure}[H]
\centering
\includegraphics[scale=0.4]{details_training}
\caption{Скриншот процесса обучения модели деталей}
\end{figure}

\begin{figure}[H]
\centering
\includegraphics[scale=0.4]{damages_training}
\caption{Скриншот процесса обучения модели повреждений}
\end{figure}

По полученным графикам можно сделать вывод о том, что модели демонстрируют сходимость в ходе обучения (снижение loss - усвоение закономерностей). Не менее показательными при этом являются не только значения функции потерь, но и итоговые значения precision, recall и mAP на валидационной и тестовой выборках (таблицы 7 и 8) - именно они показывают пригодность моделей к задаче.

Если сравнивать результаты двух моделей, то модель сегментации деталей демонстрирует результаты лучше, чем модель сегментации повреждений. Это объясняется рядом причин:
\begin{itemize}
    \item Автомобильные детали намного крупнее и имеют в большинстве случаев схожую форму
    \item Повреждения более вариативны по форме, размеру, степени выраженности и условиям съёмки
\end{itemize}

По этой причине задача сегментации повреждений является намного более сложной, а метрики по ней закономерно ниже.

С практической точки зрения сервис сохраняет прикладную ценность даже в случае неполного или неидеального распознавания - пользователь имеет возможность отредактировать результат и получить расчёт стоимости. Таким образом, ML-модели позволяют сократить объём ручной работы, но не выступают как полностью необратимая система принятия решений.

Важным результатом является визуальная интерпретируемость работы моделей. В конце работы пользователю возвращается аннотированное изображение с наложенными масками повреждений (рисунок 28), что позволяет оценить корректность работы модели. Это повышает доверие к сервису и облегчает проверку результата человеком.

На рисунках 26-30 представлены результаты работы моделей деталей и повреждений - аннотированные изображения.

\begin{figure}[H]
\centering
\includegraphics[scale=0.35]{details_inference_example}
\caption{Пример сегментации деталей на исходной фотографии}
\end{figure}

\begin{figure}[H]
\centering
\includegraphics[scale=0.4]{damages_inference_example}
\caption{Пример сегментации повреждений на обрезанных фрагментах деталей}
\end{figure}

\begin{figure}[H]
\centering
\includegraphics[scale=0.4]{annotated_user_image}
\caption{Пример итогового аннотированного изображения, отправляемого пользователю}
\end{figure}

\begin{figure}[H]
\centering
\includegraphics[scale=0.5]{many_damages_inference}
\caption{Пример корректного определения нескольких типов повреждений}
\end{figure}

\begin{figure}[H]
\centering
\includegraphics[scale=0.3]{damages_not_nice_inference}
\caption{Пример частично ошибочного определения, демонстрирующий ограничения модели}
\end{figure}

Анализируя полученные результаты, можно выделить ряд причин, по которым результат некоторых предсказаний может не соответствовать реальным повреждениям:
\begin{enumerate}[label=\arabic*.]
    \item Ограничен объём датасета и размеченных фотографий
    \item Дисбаланс классов - есть типы повреждений, которые встречаются намного реже частых, таких как царапины
    \item Сложность самих визуальных объектов, в особенности мелких царапин или вмятин
    \item Неоднозначность границ повреждений при разметке
\end{enumerate}

Достигнутый результат по ML-части можно оценить как достаточный для задачи предварительной автоматизированной оценки, но требует дальнейшего улучшения за счет увеличения кол-ва различных размеченных фотографий, в особенности, повреждений.




\section{Результаты нагрузочного тестирования}

\begin{table}[h!]
\centering
\resizebox{1.05\textwidth}{!}{
\begin{tabular}{|c|c|c|c|c|c|}
\hline
\textbf{Целевая нагрузка} & \textbf{Успешность} & \textbf{Наблюдаемый RPS} & \textbf{p50 полного цикла} & \textbf{p95 полного цикла} & \textbf{Ошибки} \\
\hline
50 RPS  & 100\%  & 37.04 RPS & 1.58 с  & 2.12 с  & - \\
\hline
120 RPS & 100\%  & 33.28 RPS & 4.40 с  & 5.62 с  & - \\
\hline
200 RPS & 93.34\% & 21.77 RPS & 12.32 с & 14.93 с & 130x ReadError \\
\hline
\end{tabular}
}
\caption{Результаты нагрузочного тестирования}
\end{table}

Нагрузочное тестирование проводилось для полного ML-сценария: создание заявки, загрузка изображения в S3-хранилище и постановка задачи инференса в Kafka. Его результаты представлены в таблице 9. При нагрузках 50 и 120 RPS обработка всех запросов прошла без ошибок. Это подтверждает устойчивость системы при повышенной нагрузке: все микросервисы работали штатно на протяжении тестирования.

На текущем одноузловом стенде система упирается примерно в 33–37 RPS. При попытке подать нагрузку выше этого значения система не выдаёт ошибок, но растёт задержка: p95 увеличивается с 2.12 с до 5.62 с.

При 200 RPS начинается деградация: успешность запросов снижается до 93.34\%, а p95 достигает почти 15 секунд. Это показывает практический предел текущего локального стенда. Основная причина ограничения - запуск backend в одном uvicorn-процессе и размещение всех сервисов на одной машине.

Итог: система выдерживает стабильную обработку около 35 RPS без ошибок, корректно буферизует ML-задачи через Kafka и показывает устойчивость при умеренной перегрузке. Для достижения уровня выше 100 RPS в продуктиве требуется горизонтальное масштабирование backend-сервиса и ML worker'ов.




\section{Результаты работы сервиса}

Пользователь взаимодействует с сервисом через бот в мессенджере Вконтакте. По итогам реализации сервис поддерживает 2 режима работы - автоматический с анализом фотографии и ручной режим ввода повреждений. 

Пользовательский сценарий начинается с команды /start (или нажатия кнопки "Начать") с последующим выбором режима работы (Рисунок 31). При выборе режима с фотографией пользователь отправляет изображение автомобиля с видимым повреждением (Рисунок 32). 

\begin{figure}[H]
\centering
\includegraphics[scale=0.55]{start_interface}
\caption{Интерфейс бота (начальный этап выбора режима работы)}
\end{figure}

\begin{figure}[H]
\centering
\includegraphics[scale=0.75]{ml_mode_interface}
\caption{Пример взаимодействия с ботом в режиме ML}
\end{figure}

Сервис сохраняет фотографию, после чего система выполняет обработку изображения и возвращает аннотированное изображение вместе со списком найденных повреждений и клавиатуру с возможностью отредактировать повреждения и кнопкой подтверждения (Рисунок 33).

\begin{figure}[H]
\centering
\includegraphics[scale=0.8]{ml_mode_result}
\caption{Результат обработки фотографии в режиме ML}
\end{figure}

Существенным преимуществом в интерфейсе бота является функционал, где при необходимости пользователь может отредактировать список обнаруженных повреждений - удалить, изменить тип найденного повреждения или его принадлежность к детали, или добавить необнаруженное повреждение (Рисунки 34 и 35). Это делает систему более пригодной для практического пользования и не делает пользователя зависимым от возможных просчётов модели.

\begin{figure}[H]
\centering
\includegraphics[scale=0.85]{delete_damage}
\caption{Редактирование списка повреждений (удаление)}
\end{figure}

\begin{figure}[H]
\centering
\includegraphics[scale=0.85]{interface_damages}
\caption{Редактирование списка повреждений (добавление)}
\end{figure}

После того, как пользователь введёт полный список повреждений, он может нажать кнопку "Подтвердить", и сервис сразу выдаст полную стоимость и длительность работ в рабочих часах (Рисунок 36):

\begin{figure}[H]
\centering
\includegraphics[scale=0.85]{price_result}
\caption{Результат подсчёта стоимости и длительности работ}
\end{figure}

Пользователь может выбрать и режим ручного ввода повреждений. В этом режиме ему доступен схожий функционал, что и в режиме ML при редактировании повреждений, только теперь пользователь вводит все повреждения с нуля самостоятельно. (Рисунок 37)

\begin{figure}[H]
\centering
\includegraphics[scale=0.8]{manual_mode}
\caption{Пример взаимодействия с ботом в ручном режиме}
\end{figure}

Таким образом, с точки зрения пользовательского опыта достигнут результат, при котором:
\begin{itemize}
    \item Сценарий запуска является быстрым и понятным
    \item Доступна свобода действий и выбора автоматического и ручного режима
    \item Пользователь имеет удобный, не перегруженный интерфейс и может контролировать результат
    \item Итоговая смета представляется в понятном виде
\end{itemize}

\section{Оценка выполнения задач исследования}

В рамках работы был подготовлен датасет изображений деталей и повреждений и выполнена разметка в формате, пригодном для дообучения моделей YOLOv8-seg. Уровень решения задачи можно оценить как высокий с точки зрения работоспособного датасета и средний с точки зрения охвата всех редких случаев. Ограничением всё так же остаётся необходимости дальнейшего расширения и балансировки выборки, особенно для сложных и трудноразличимых повреждений.

Разработка модулей для обработки изображений также решена. На этапе обучения используется конфигурируемый набор аугментаций, что позволило модели быть устойчивой к разнообразию фотографий. Использование аугментаций также помогло частично компенсировать ограниченный объём размеченных данных. 
Недостатки съёмки, такие как нечёткие фотографии, блики, плохое освещение, всё же существенно снижают качество распознавания, так как учесть все возможные вариации съёмки и передать эту информацию модели невозможно. 
Однако обучение моделей показало, что они демонстрируют сходимость в ходе обучения, что видно по функции потерь и метрикам - это означает, что модели пригодны к задаче, и главной проблемой является необходимость большого количества грамотно размеченных данных. 

Бизнес-логика оценки стоимости и длительности ремонта реализована в полном объёме. Система корректно учитывает все особенности предметной области и обладает знанием о всех расценках автосервиса и корректно применяет их ко всем парам деталей и повреждений, как это делают сотрудники автосервиса. 

Система построена в виде микросервисной архитектуры, включающей Backend API, Bot Service и ML Worker. Обмен сообщениями осуществляется через Kafka, для хранения изображений используется S3, а данные заявок и повреждений сохраняются в PostgreSQL.  Такое архитектурное решение позволило разделить ответственность между разными компонентами системы, сделать её масштабируемой. Интеграция с ботом в мессенджере выполнена - пользователь может полностью пройти весь сценарий оценки, не покидая интерфейс Вконтакте.
Архитектура базы данных соответствует микросервисному подходу и поддерживает основные сценарии работы системы.

Также произведена интеграция CI/CD - настроен конвейер с использованием GitHub Actions, выполняющий автоматическую проверку стиля кода и запуск тестов при изменениях в целевой ветви. Система контейнеризирована с использованием Docker Compose, что обеспечивает воспроизводимую локальную сборку, и развёрнута на сервере.

Проведено полноценное тестирование работоспособности системы - в проекте реализовано более 500 юнит- и интеграционных тестов, которые поддерживают устойчивость реализованных сценариев на уровне кода. 
Дополнительно реализован отдельный скрипт нагрузочного тестирования, имитирующий полный сценарий ML режима работы: создание заявки, загрузка изображений в объектное хранилище. Это позволило оценить пропускную способность системы и устойчивость к интенсивному потоку запросов. Результаты нагрузочного тестирования были представлены выше в таблице 9.

На основании приведённого анализа можно сделать вывод о том, что все поставленные задачи получили практическое решение. Наибольший потенциал для дальнейшего улучшения остаётся в повышении качества ML-моделей и разметки большего количества фотографий деталей и повреждений.

Следует отметить, что достигнутый результат не означает полной замены эксперта-оценщика. Итоговая система является инструментом предварительной автоматизированной оценки, а не окончательным юридически значимым подсчётом стоимости ремонта.

